
\section{Lean formalization plan}
\label{sec:lean}

\subsection{Target namespace and library placement}

The implementation should become a new sublibrary under
\[
 \code{Logic/PrimitiveRecursiveArithmetic/Derivations/Lean/IPRADerivations/Framework}.
\]
The proposed namespace is
\[
       \code{LeanProofs.IPRA.Framework}.
\]
It should depend on \code{ConciseIPRA.PR} and selected checked
\code{IPRADerivations.Core} files.  It should not depend on theorem packets.

A first Lake target may import:
\[
\begin{array}{l}
\code{IPRADerivations.Core.Basic},\\
\code{IPRADerivations.Core.Definitions},\\
\code{IPRADerivations.Core.Order},\\
\code{IPRADerivations.Core.Pairing},\\
\code{IPRADerivations.Core.Search}.
\end{array}
\]
Later files may import bounded sums, rational coding, matrices, or bridge
modules as needed.

\subsection{Phase 0: structural wrapper}

The first file is almost entirely checked reuse.

\begin{lstlisting}[language={}]
structure Module (n : Nat) where
  req   : List (Fm PTm n)
  out   : Fm PTm n
  proof : req |- [PRAAx] out

namespace Module

def specialize (M : Module n) (sigma : Fin n -> PTm m) :
    Module m where
  req := M.req.map (fun p => p.subst sigma)
  out := M.out.subst sigma
  proof := M.proof.subst sigma

def discharge (M : Module n) {Delta : List (Fm PTm n)}
    (h : forall p, List.Mem p M.req -> Delta |- [PRAAx] p) :
    Delta |- [PRAAx] M.out :=
  M.proof.trans h
\end{lstlisting}

The actual source uses the repository's scoped turnstile notation.  The archive
contains an uncompiled candidate spelling.  The acceptance test is:
\[
    \code{lake build IPRADerivations}.
\]

Add small theorems for:
\begin{itemize}
\item weakening a module context;
\item turning a closed proof into a module with empty requirements;
\item a unary \code{Contract};
\item contract identity and composition;
\item conjunction introduction/elimination modules;
\item context permutation/union conveniences.
\end{itemize}

\subsection{Phase 1: Boolean core}

Create:
\[
\begin{array}{l}
\code{Framework/Boolean.lean},\\
\code{Framework/ReflectQF.lean},\\
\code{Framework/BoundedQuantifier.lean}.
\end{array}
\]

\code{Boolean.lean} defines term operations \code{bnot}, \code{band},
\code{bor}, and \code{bimp}, proves substitution lemmas, Boolean closure, and
truth tables.  Reuse \code{eqb}, \code{eqb_of_eq}, and \code{eqb_of_ne} from
\code{Core.Pairing}.

The central structure should be close to:

\begin{lstlisting}[language={}]
structure ReflectedQF {n : Nat} (phi : Fm PTm n) where
  bit       : PTm n
  isBool    : ([] : List (Fm PTm n)) |- [PRAAx]
                ((bit = 0) OR (bit = 1))
  trueIff   : ([] : List (Fm PTm n)) |- [PRAAx]
                iff phi (bit = 1)
  falseIff  : ([] : List (Fm PTm n)) |- [PRAAx]
                iff (phi --> bot) (bit = 0)
\end{lstlisting}

Use repository notation in the final source.  Implement:

\begin{lstlisting}[language={}]
def reflectQF (phi : Fm PTm n) (hphi : phi.QF) :
    ReflectedQF phi
\end{lstlisting}

by recursion on \(\varphi\), with impossible quantifier cases eliminated by
\code{hphi}.  A useful auxiliary theorem is internal decidability:

\begin{lstlisting}[language={}]
theorem qf_decidable (r : ReflectedQF phi) :
    Gamma |- [PRAAx] phi OR (phi --> bot)
\end{lstlisting}

Then implement arbitrary-arity finite folds by encoding the parameter tuple
with pairing, or generalize \code{recDef} directly.  The initial version can
use one packed parameter because the existing \code{Search} API is binary.

Acceptance tests:
\begin{itemize}
\item reflection of equality, conjunction, disjunction, implication, and
falsity;
\item substitution commutation on a formula with nested connectives;
\item reflection of the order formulas already defined in \code{Core.Order};
\item no new axioms and no \code{sorry}.
\end{itemize}

\subsection{Phase 2: trace invariant package}

Create \code{Framework/Trace.lean}.  Rather than building a deep program
syntax, define a record of the already specialized ingredients.

A possible record is:

\begin{lstlisting}[language={}]
structure TraceSpec (k : Nat) where
  base : PTm k
  step : PTm (k + 2)
  inv  : Fm PTm (k + 1)
  invQF : inv.QF
  baseProof :
    ([] : List (Fm PTm k)) |- [PRAAx]
      inv.subst (inst base)
  stepProof :
    [inv] |- [PRAAx]
      inv.subst (repl (step.subst ...))
\end{lstlisting}

The indices in the sketch must be chosen to match the convention:
variable \(0\) is the current state, variable \(1\) may be time, and the
remaining variables are parameters.  It may be cleaner to store the fully
formed induction formula \code{runInv : Fm PTm (k+1)} and its base/step
derivations, then define a higher-level constructor that produces these from
\code{base}, \code{step}, and \code{inv}.  This separates binder bookkeeping
from the mathematical theorem.

The endpoint theorem should be:

\begin{lstlisting}[language={}]
theorem TraceSpec.invariant (S : TraceSpec k) :
    ([] : List (Fm PTm (k + 1))) |- [PRAAx] S.runInv
\end{lstlisting}

implemented by \code{IPRA.ind}.  Add a terminal contract combinator that
specializes and cuts an invariant proof with a halt proof.

Acceptance example: reprove a small checked bounded fold, not a theorem packet.
For instance, prove the invariant of finite sum or a running maximum.  This
tests substitution and state packing before a more complicated application.

\subsection{Phase 3: bounded relations and stability}

Create:
\[
\begin{array}{l}
\code{Framework/Relation.lean},\\
\code{Framework/Viability.lean},\\
\code{Framework/Observer.lean}.
\end{array}
\]

A bounded relation package stores a reflected predicate, candidate bound, and
possibly a totality module.  Backward viability is a primitive recursion over
time whose state is a bit table or a packed sequence.  For a first
implementation, avoid a general array library by using a characteristic
function encoded through a bounded table code with checked lookup/update
equations.

Required theorems:
\[
\begin{array}{ll}
\code{viable_succ}:&
 \Viable_t(a)=1\to
 \exists b<B_{t+1}\,[R_t(a,b)=1\wedge\Viable_{t+1}(b)=1],\\
\code{selectPath_spec}:&
 \Viable_0(a_0)=1\to \mathsf{PathGood}(\mathsf{selectPath}(\cdots))=1,\\
\code{stable_exact}:&
 \text{pairwise equality}\to \text{selection-independent output},\\
\code{stable_setoid}:&
 \text{pairwise output equivalence}\to \text{extensional result},\\
\code{stable_error}:&
 \text{diameter bound}\to \text{observer error bound}.
\end{array}
\]

The exact observer theorem should be short because it is mainly
\code{Search.search_found}, hypothesis specialization, and equality
substitution.  Implement it before the full viability table.

\subsection{Phase 4: graph assembler}

The first graph assembler may stay outside the object theory.  Define a Lean
inductive type whose leaves contain \code{Module} values and whose internal
nodes contain functions consuming predecessor proofs.  Alternatively, avoid
an inductive syntax entirely and expose a combinator API.

A minimal useful API includes:
\[
\begin{array}{l}
\code{Module.specialize},\quad
\code{Module.discharge},\quad
\code{Module.and},\quad
\code{Module.branch},\\
\code{Module.loop},\quad
\code{Module.search},\quad
\code{Module.observe},\quad
\code{Module.residual}.
\end{array}
\]

A deep serialized graph can be added later.  Its elaborator must return
\code{Module n}; that result type is the soundness boundary.

Add a metadata structure outside the trusted proof:
\begin{lstlisting}[language={}]
structure ObligationInfo where
  name   : String
  kind   : ObligationKind
  source : String
  note   : String
\end{lstlisting}
A compiled artifact contains both the proof module and an obligation ledger.
The ledger may be JSON for repository tooling.

\subsection{Phase 5: open bridge wrappers}

The bridge work should be split into small files.

\paragraph{\code{CodeNormalization.lean}.}
Prove extensional equality of each Foundation construction code with its
Foundation--Concise--Foundation normal form.  The primitive-recursion case is
the only induction.

\paragraph{\code{FormulaRoundtrip.lean}.}
Lift code equality to terms and formulas, first equality-only, then native
order by composing with \code{nativeOrderRoundtrip}.

\paragraph{\code{OpenSequent.lean}.}
Package:
\begin{lstlisting}[language={}]
theorem conciseToStandardCtx ...
theorem standardToConciseCtx ...
theorem standard_concise_ctx_iff ...
\end{lstlisting}
with both arbitrary environments and a canonical variable environment.

The existing \code{toFoundation_provable_of_axioms} should be reused directly.
For the forward direction, initially use finite-context internalization,
universal closure, the current closed translator, repeated
\code{Prf.allOpen}, and implication elimination.  If performance is poor, add
a direct context-carrying version of \code{hilbertToConcise}.

\subsection{Phase 6: domain libraries}

Only after the core compiles should domain presentations be added:

\[
\begin{array}{ll}
\code{Framework/Data/Setoid.lean} & generic validity/equivalence interfaces,\\
\code{Framework/Data/Rational.lean} & signed rationals and inequalities,\\
\code{Framework/Data/FiniteSet.lean} & list setoids and deduplication,\\
\code{Framework/Analysis/RegularName.lean} & rational names and moduli,\\
\code{Framework/Topology/FiniteNerve.lean} & finite covers and nerves,\\
\code{Framework/Algebra/PolynomialCert.lean} & identity and ideal certificates.
\end{array}
\]

Existing theorem packets should migrate by importing these libraries rather
than defining isolated encodings.

\subsection{Proof order and dependency discipline}

The recommended order is:

\begin{enumerate}
\item structural \code{Module};
\item Boolean term truth tables;
\item QF reflection;
\item bounded universal/existential reflection;
\item trace invariant;
\item exact stable observer;
\item finite viability;
\item approximate observer;
\item open bridge normalization;
\item domain presentations;
\item theorem-packet migration.
\end{enumerate}

Each step has a standalone test theorem.  Do not begin with a theorem packet
that depends simultaneously on rational analysis, finite arrays, bridge
transport, and relational stability; that would obscure which layer failed.

\subsection{No-placeholder rule}

A file is placed in the checked library only when:
\begin{itemize}
\item it builds under the pinned Mathlib and repository commit;
\item it contains no \code{sorry}, \code{admit}, new axiom, or opaque semantic
oracle;
\item every public theorem concludes a \code{Prf PRAAx} judgment or a clearly
labeled metatheorem used to construct one;
\item its README distinguishes object-level theorems from host-level tests.
\end{itemize}

Blueprint files may live in a proposal directory, but they should not be
counted as checked derivations.

\subsection{Current artifact status}

Lean is not installed in the execution environment used to prepare this
archive.  The included \code{lean/} files are therefore source-oriented
skeletons and exact signature proposals, not a claimed successful build.
They intentionally avoid presenting placeholders as completed proofs.  The
first repository-side action is to adapt \code{Core.lean} to exact notation
and run the build target; subsequent files can then be filled in one theorem
at a time.
